% =====================================================================
%  CAPÍTULO 6 -- RESULTADOS E CONSIDERAÇÕES FINAIS
%  (Regulamento, Art. 9º, item vi; Manual, seção 2.2.3)
%  O Manual recomenda que a parte conclusiva NÃO contenha citações.
% =====================================================================

\chapter{Resultados e considerações finais}
\label{cap:resultados}

Este projeto conduziu a concepção do RoundUp como um processo de
ideação do autor, no qual as ideias de funcionalidades foram convertidas
em requisitos, construídas como hipóteses e, em paralelo, submetidas a um
questionário que as confirmou ou refutou e revelou ideias novas. Essa
estratégia permitiu rastrear cada funcionalidade desde a ideia até a
decisão de mantê-la, ajustá-la ou revisá-la, e daí até o código e os
testes.

\section{Resultados alcançados}

Os objetivos específicos apresentados no \autoref{cap:introducao} foram
atingidos da seguinte forma:

\begin{enumerate}
  \item \textbf{Fatores de motivação e evasão.} Atingido de forma
        exploratória e limitada. Os praticantes reais indicaram a saúde
        mental, a saúde física e o convívio com a turma como motivos para
        continuar treinando, e os dois ex-praticantes pararam por falta de
        tempo ou conflito com a rotina. Ambos afirmaram que um aplicativo
        não teria mudado sua decisão, o que sugere que parte da evasão tem
        causa externa ao treino. A visão dos professores sobre a evasão não
        pôde ser levantada.
  \item \textbf{Ideias convertidas em requisitos rastreáveis.} Atingido.
        Treze ideias originaram, junto com os requisitos de base, 33
        requisitos funcionais, 9 regras de negócio e 11 requisitos não
        funcionais, todos registrados com origem, ciclo, prioridade e
        situação.
  \item \textbf{Confirmação ou refutação com usuários reais.} Atingido,
        com amostra pequena. Seis respondentes avaliaram as ideias: seis
        foram confirmadas, com destaque para os recursos de progresso
        visível, e seis não foram, entre elas os lembretes, a mensagem
        automática após ausência, as missões e o próprio
        \textit{check-in}, que segue implementado por ser a base do
        sistema. Os comentários confirmaram riscos que o projeto já
        tratava, como o esforço de registro e o excesso de notificações,
        e trouxeram duas ideias novas.
  \item \textbf{Personas: ideias novas, sinal complementar e teste do
        instrumento.} Atingido. Doze personas sintéticas orientaram a
        geração de 65 respostas, enviadas a uma cópia do formulário. Como
        fonte de ideias, elas trouxeram sete dos nove requisitos
        candidatos registrados, quatro deles do ponto de vista do
        professor, que a coleta real não alcançou. Como sinal
        complementar, convergiram com os dados reais na maior parte das
        ideias e divergiram no \textit{ranking} e nas missões, apontando
        onde investigar. Como teste, confirmaram o roteamento das três
        seções e o funcionamento do procedimento de análise, que fica
        pronto para uma nova coleta real. Em nenhum caso decidiram sobre
        um requisito.
  \item \textbf{Modelagem e implementação do protótipo com apoio de
        agentes de IA.} Atingido. O protótipo é um aplicativo funcional
        com três perfis de usuário. Os 33 requisitos funcionais foram
        implementados; a interface foi conferida na versão web, e a data
        de conquista das medalhas, o público-alvo e a vigência das
        missões, a personalização das faixas e o envio de notificações
        foram verificados em um iPhone. As regras de negócio que protegem
        a gamificação, como a validação humana antes do crédito de pontos,
        foram implementadas nas regras de segurança do banco de dados e
        cobertas por 288 testes automatizados. A maior parte do código foi
        escrita por um agente de IA, a partir de requisitos, roteiros e
        planos do autor, que concentrou o seu trabalho no planejamento, no
        projeto e no controle de qualidade.
  \item \textbf{Conformidade com a LGPD.} Atingido. O consentimento no
        cadastro, o registro do responsável legal de alunos menores de
        idade, a guarda dos dados pessoais em documento visível só ao
        próprio usuário, o isolamento dos dados entre academias e professores e a
        exclusão de conta com remoção dos dados em cascata foram
        implementados. Fica em aberto se o registro do consentimento deve
        ser preservado após a exclusão, como prova, o que depende de
        orientação jurídica.
\end{enumerate}

\section{Considerações sobre o uso de respostas sintéticas}

As personas e as respostas sintéticas foram úteis para dois fins: como
insumo complementar da ideação, ao dar voz a perfis que a coleta real não
alcançou e levar o autor a repensar os requisitos, e como dados de teste
do questionário e da análise, principalmente na seção do professor. Elas não substituem a coleta com
usuários reais, porque reproduzem padrões supostos por quem as gerou;
quando ficam próximas dos dados reais, isso mostra apenas que o gerador
copiou bem o padrão que lhe deu origem. A lição prática é manter os dois
conjuntos sempre separados, identificar cada resposta sintética e
registrar o método de geração, para que nenhuma conclusão sobre os
usuários dependa delas.

\section{Considerações sobre o desenvolvimento assistido por IA}

O uso de um agente de codificação mudou o que coube ao autor fazer, mas
não reduziu a sua responsabilidade. O agente escreveu código, testes e
registros de decisão com rapidez, mas foram os artefatos do autor que
mantiveram o trabalho segundo princípios básicos de engenharia e qualidade: a lista de requisitos com códigos
rastreáveis, os roteiros com critérios de aceite, a exigência de um plano
antes de cada mudança grande e, sobretudo, as regras de segurança e os
testes, que verificaram o que foi gerado de forma independente. Os
defeitos que só apareceram nos testes e na execução em aparelho mostram
por que essa verificação não é opcional. A lição prática é que, com
agentes de IA, o valor do desenvolvedor se desloca da digitação do código
para a capacidade de especificar bem, projetar a solução, dividir o
trabalho em etapas pequenas e verificar cada uma delas. É essa
capacidade, mais do que a ferramenta, que determina a qualidade do
resultado.

\section{Limitações}

\begin{itemize}
  \item A amostra real é pequena (seis respondentes), pouco diversa
        (todos de Jiu-Jitsu) e não incluiu professores; seus resultados
        servem para orientar decisões, não para generalizar. No grupo de
        ex-praticantes, com duas pessoas, as médias mais baixas dependem
        das notas de um único respondente.
  \item A avaliação com pessoas incidiu sobre as ideias, apresentadas
        como descrições de recursos; o aplicativo construído ainda não foi
        usado por alunos ou professores reais, e foi verificado em aparelho
        apenas em iPhone.
  \item Nenhum professor real respondeu, de modo que os requisitos de
        gestão não foram avaliados; as respostas sintéticas de professor
        servem apenas como dados de teste.
  \item Os valores-padrão de gamificação (pontos por presença, níveis e
        marcos de medalha) não foram calibrados com uso real, e o lembrete
        diário agendado das notificações foi testado apenas em emulador.
  \item O efeito do aplicativo sobre a evasão não pôde ser medido no prazo
        do TCC, e parte da evasão tem causas externas, como falta de tempo
        e lesões, que estão fora do alcance de um aplicativo.
  \item O histórico do repositório não distingue o código escrito pelo
        agente de IA do escrito pelo autor, e não houve condição de
        comparação sem IA; por isso, o relatório descreve o processo de
        desenvolvimento assistido, mas não mede o seu ganho de
        produtividade.
\end{itemize}

\section{Trabalhos futuros}

\begin{itemize}
      \item Validar em aparelho o lembrete diário agendado e fazer o toque
        na notificação levar à tela correspondente (RF31).
  \item Permitir o uso sem conexão (RNF06) nas ações que não dependem de
        validação imediata do servidor, mantendo o \textit{check-in} e a
        validação de presença dependentes de conexão.
  \item Reaplicar o questionário com uma amostra maior, incluindo
        outras modalidades, concentrando a investigação nas ideias marcadas
        como ``a revisar''.
  \item Aplicar o questionário a professores e donos de academia reais,
        para avaliar os requisitos de gestão hoje marcados como ``sem
        dados''.
  \item Realizar um teste de uso do aplicativo com alunos e professores
        de uma academia, incluindo a execução em aparelhos Android.
  \item Investigar os requisitos candidatos do \quadroref{qua:ideias-novas},
        começando pelos de origem sintética, que ainda precisam ser
        confirmados por pessoas reais.
  \item Realizar um estudo de acompanhamento, medindo a frequência dos
        alunos antes e depois do uso do aplicativo.
\end{itemize}